\scanpage{15}
\mspage{14}\mspage{7}
\margin{on suppose $M5$ mais pas $M6$ ni $M7$\ldots}
\figure{\begin{tikzpicture}\node[draw,ellipse] {$\overline{\underline{\mathrm{IV}}}$};\end{tikzpicture}}
Ind-structures modérées strictes \inserted{\uline{indespaces $\overset{\overline{\longrightarrow}}{\text{modérés}}$}} \inserted{(ou ind-espaces modérés)} : \unsure{consiste} : la donnée d'un ens. $E$, et un ens de parties $\mathcal M_E$ sur $E$ (les parties modérées), stable par réunions finies, \inserted{de réunion $E$}, avec une structure modérée sur chacun, ces structures s'induisant mutuellement, et
\[
 A\in\mathcal M,\quad B\in\mathcal M_A\Longrightarrow B\in\mathcal M.
\]
Il revient au $=$ de se donner la notion d'application modérée d'un $X$ modéré dans $E$, avec \deleted{une identité} quelques conditions supplém.. Si on suppose
\[
 (M8)\quad R\text{ est réunion \inserted{(filtrante par $M5$)} de ses parties modérées},
\]
on trouve une structure ind-modérée canonique sur $R$, et de $=$ sur les $R^I$.

\uline{Proposition} La catégorie des \inserted{$\overset{\leftrightarrow}{\underrightarrow{\operatorname{Esp}}}_{\mathcal M_*}$} \deleted{Ind} Ind-espaces modérés stricts admet des $\varprojlim$ finies, le foncteur oubli y commute.
\margin{\unsure{des ensembles modérés|des ind-modérés}}
\margin{$M1:M5$}
\margin{$\underrightarrow{\operatorname{Esp}}_{\mathcal M_*}$ est obtenue par $\uline{\varinjlim}$ filtrantes \unsure{des injections} \deleted{strictes} \unsure{dans les ensembles} \unsure{modérés} \unsure{quelconques}.}

Si $R_0\ne\varnothing$ (idem dans $\underrightarrow{\operatorname{Esp}}_{\mathcal M_*}$ il y a sommes disjointes finies) il y a des sommes disjointes \deleted{finies} \inserted{(en principe !)} dans $\overset{\leftrightarrow}{\underrightarrow{\operatorname{Esp}}}_{\mathcal M_*}$, et $X\longrightarrow |X|$ y commute.
